%-------------------------
% Resume in Latex
% Author : Jake Gutierrez
% Based off of: https://github.com/sb2nov/resume
% License : MIT
%------------------------

\documentclass[letterpaper,11pt]{article}

\usepackage{latexsym}
\usepackage[empty]{fullpage}
\usepackage{titlesec}
\usepackage{marvosym}
\usepackage[usenames,dvipsnames]{color}
\usepackage{verbatim}
\usepackage{enumitem}
\usepackage[hidelinks]{hyperref}
\usepackage{fancyhdr}
\usepackage[english]{babel}
\usepackage{tabularx}
%\usepackage{cite}
\usepackage{csquotes}
\usepackage{xcolor}
\definecolor{papertitlelink}{HTML}{0066CC}
\input{glyphtounicode}


%\usepackage[style=ieee,url=false,doi=false,maxbibnames=99,sorting=ydnt,dashed=false,backend=biber]{biblatex}
\usepackage[style=ieee,url=false,doi=false,maxbibnames=99,sorting=none,dashed=false,citestyle=numeric-comp,backend=biber]{biblatex}
% html という独自フィールドを biblatex/biber に認識させる
% bib の html フィールドを biblatex 標準フィールド usera に写す
\DeclareSourcemap{
  \maps[datatype=bibtex]{
    \map{
      \step[fieldsource=html, fieldtarget=usera]
    }
    \map{
      \step[fieldsource=arxiv, fieldtarget=userb]
    }
    \map{
      \step[fieldsource=equal_contribution, fieldtarget=userc]
    }
  }
}

% タイトルリンク用の共通マクロ
% 優先順位: html > arxiv > リンクなし
\newcommand{\linkedbibtitle}[1]{%
  \iffieldundef{usera}
    {%
      \iffieldundef{userb}
        {#1}
        {\href{https://arxiv.org/abs/\thefield{userb}}{\textcolor{papertitlelink}{#1}}}%
    }
    {\href{\thefield{usera}}{\textcolor{papertitlelink}{#1}}}%
}
% inproceedings に限らず title フィールドすべてに適用
\DeclareFieldFormat*{title}{%
  \mkbibquote{\linkedbibtitle{#1\isdot}}%
}

\addbibresource{journal.bib}
\addbibresource{conference.bib}
\addbibresource{preprint.bib}

\usepackage{customcv}
\boldname{Horiguchi}{Shota}{S.}
\AtBeginBibliography{\small}


%----------FONT OPTIONS----------
% sans-serif
% \usepackage[sfdefault]{FiraSans}
% \usepackage[sfdefault]{roboto}
% \usepackage[sfdefault]{noto-sans}
% \usepackage[default]{sourcesanspro}

% serif
% \usepackage{CormorantGaramond}
% \usepackage{charter}


\pagestyle{fancy}
\fancyhf{} % clear all header and footer fields
\fancyfoot{}
\renewcommand{\headrulewidth}{0pt}
\renewcommand{\footrulewidth}{0pt}

% Adjust margins
\addtolength{\oddsidemargin}{-0.5in}
\addtolength{\evensidemargin}{-0.5in}
\addtolength{\textwidth}{1in}
\addtolength{\topmargin}{-.5in}
\addtolength{\textheight}{1.0in}

\urlstyle{same}

\raggedbottom
\raggedright
\setlength{\tabcolsep}{0in}

% Sections formatting
\titleformat{\section}{
  \vspace{-4pt}\scshape\raggedright\large
}{}{0em}{}[\color{black}\titlerule \vspace{-5pt}]

% Ensure that generate pdf is machine readable/ATS parsable
\pdfgentounicode=1

%-------------------------
% Custom commands
\newcommand{\resumeItem}[1]{
  \item\small{
    {#1 \vspace{-2pt}}
  }
}
\makeatletter
\newcommand{\resumeSubheading}[2]{
  \vspace{-2pt}\item
    \begin{tabular*}{0.97\textwidth}[t]{l@{\extracolsep{\fill}}r}
      \textbf{#1} & #2
      %\textit{\small#3} & \textit{\small #4} \\
      \@resumeSubheadingIter
}
\newcommand\@resumeSubheadingIter{% 
    \@ifnextchar\bgroup{\@resumeSubheadingPosition}{\@resumeSubheadingEnd}}
\newcommand\@resumeSubheadingPosition[2]{
    \\\textit{\small #1} & \textit{\small #2}
    \@resumeSubheadingIter
}%&\textit{#2}

\newcommand\@resumeSubheadingEnd{
    \end{tabular*}\vspace{-7pt}
}
\makeatother

\newcommand{\resumeShortSubheading}[2]{
  \vspace{-2pt}\item
    \begin{tabular*}{0.97\textwidth}[t]{l@{\extracolsep{\fill}}r}
      \textbf{#1} & #2 \\
    \end{tabular*}\vspace{-7pt}
}

\newcommand{\resumeSubSubheading}[2]{
    \item
    \begin{tabular*}{0.97\textwidth}{l@{\extracolsep{\fill}}r}
      \textit{\small#1} & \textit{\small #2} \\
    \end{tabular*}\vspace{-7pt}
}

\newcommand{\resumeProjectHeading}[2]{
    \item
    \begin{tabular*}{0.97\textwidth}{l@{\extracolsep{\fill}}r}
      \small#1 & #2 \\
    \end{tabular*}\vspace{-7pt}
}

\newcommand{\resumeSubItem}[1]{\resumeItem{#1}\vspace{-4pt}}

\renewcommand\labelitemii{$\vcenter{\hbox{\tiny$\bullet$}}$}

\newcommand{\resumeSubHeadingListStart}{\begin{itemize}[leftmargin=0.15in, label={}]}
\newcommand{\resumeSubHeadingListEnd}{\end{itemize}}
\newcommand{\resumeItemListStart}{\begin{itemize}}
\newcommand{\resumeItemListEnd}{\end{itemize}\vspace{-5pt}}

%-------------------------------------------
%%%%%%  RESUME STARTS HERE  %%%%%%%%%%%%%%%%%%%%%%%%%%%%


\begin{document}
\nocite{*}
%----------HEADING----------
% \begin{tabular*}{\textwidth}{l@{\extracolsep{\fill}}r}
%   \textbf{\href{http://sourabhbajaj.com/}{\Large Sourabh Bajaj}} & Email : \href{mailto:sourabh@sourabhbajaj.com}{sourabh@sourabhbajaj.com}\\
%   \href{http://sourabhbajaj.com/}{http://www.sourabhbajaj.com} & Mobile : +1-123-456-7890 \\
% \end{tabular*}
(Updated on Jun. 5, 2026)
\begin{center}
    \textbf{\Huge \scshape Shota Horiguchi} \\ \vspace{1pt}
    \href{mailto:horiguchi@ieee.org}{\underline{horiguchi@ieee.org}} $|$ 
    \href{mailto:shota.horiguchi@ntt.com}{\underline{shota.horiguchi@ntt.com}} $|$ 
    \href{shota-horiguchi.github.io}{\underline{shota-horiguchi.github.io}} $|$
    \href{linkedin.com/in/shota-horiguchi}{\underline{linkedin.com/in/shota-horiguchi}}
    %\href{https://github.com/...}{\underline{github.com/hshota0530}}
\end{center}


%-----------EDUCATION-----------
\section{Education}
  \resumeSubHeadingListStart
    \resumeSubheading
      {University of Tsukuba}{Ibaraki, Japan}
      {Ph.D. in Computer Science}{Apr. 2022 -- Mar. 2023}
      \resumeItemListStart
        \resumeItem{Thesis title: Study on overlap-aware speaker diarization and its applications}
        \resumeItem{Supervisor: Prof. Takeshi Yamada}
      \resumeItemListEnd
    \resumeSubheading
      {The University of Tokyo}{Tokyo, Japan}
      {Master of Information Science and Technology}{Apr. 2015 -- May 2017}
      \resumeItemListStart
        \resumeItem{Thesis title: Personalized object recognition}
        \resumeItem{Supervisor: Prof. Kiyoharu Aizawa}
        \resumeItem{Related publications:~\cite{horiguchi2018personalized,horiguchi2020significance}}
      \resumeItemListEnd
    \resumeSubheading
      {The University of Tokyo}{Tokyo, Japan}
      {Bachelor of Engineering}{Apr. 2011 -- May 2015}
      \resumeItemListStart
        \resumeItem{Supervisor: Prof. Kiyoharu Aizawa}
        \resumeItem{Related publications:~\cite{horiguchi2016lognormal}}
      \resumeItemListEnd
    \resumeSubHeadingListEnd


%-----------EXPERIENCE-----------
\section{Work Experiences}
  \resumeSubHeadingListStart
    \resumeSubheading
      {NTT, Inc., Human Informatics Laboratories}{Kanagawa, Japan}
      {Research Specialist}{February 2024 -- Present}
      \resumeItemListStart
        \resumeItem{Assessing annotation ambiguity in speaker diarization (2025--Present)~\cite{horiguchi2025can,horiguchi2026tight}}
        \resumeItem{Speaker embedding extraction from multi-speaker recordings (2024--Present)~\cite{horiguchi2024recursive,horiguchi2025guided,horiguchi2025mitigating,horiguchi2025pretraining}}
      \resumeItemListEnd
    \resumeSubheading
      {Hitachi, Ltd. Research and Development Group}{Tokyo, Japan}
      {Senior Researcher}{October 2021 -- January 2024}
      {Researcher}{April 2017 -- September 2021}
      \resumeItemListStart
        \resumeItem{Maintaining machine learning models (2022--2024)~\cite{horiguchi2024streaming}}
        \resumeItem{Speaker diarization, speech separation, and meeting transcription using distributed microphones (2019--2023)~\cite{horiguchi2020utterancewise,horiguchi2021blockonline,horiguchi2022multichannel,horiguchi2023mutual}}
        \resumeItem{Speaker diarization for cases where the number of speakers is unknown (2019--2023)~\cite{horiguchi2022encoderdecoder,horiguchi2023online,horiguchi2020endtoend,horiguchi2021towards,horiguchi2021endtoend}}
        \resumeItem{Multimodal recognition and interaction for humanoid robots (2017--2019)~\cite{horiguchi2018facevoice,horiguchi2019multimodal}}
      \resumeItemListEnd
    \resumeSubheading
      {Yahoo Japan Corporation}{Tokyo, Japan}
      {Research \& Development Intern}{August, October -- December 2015}
      \resumeItemListStart
        \resumeItem{Click through rates prediction}
      \resumeItemListEnd
% -----------Multiple Positions Heading-----------
%    \resumeSubSubheading
%     {Software Engineer I}{Oct 2014 - Sep 2016}
%     \resumeItemListStart
%        \resumeItem{Apache Beam}
%          {Apache Beam is a unified model for defining both batch and streaming data-parallel processing pipelines}
%     \resumeItemListEnd
%    \resumeSubHeadingListEnd
%-------------------------------------------+
  \resumeSubHeadingListEnd


%-----------PROJECTS-----------
\begin{comment}
\section{Projects}
    \resumeSubHeadingListStart
      \resumeProjectHeading
          {\textbf{Concept drift detection and model adaptation}}{April 2022 -- January 2024}
          \resumeItemListStart
            \resumeItem{Developed a meeting transcription framework based on distributed microphones, e.g., smartphones or tablets \cite{horiguchi2020utterancewise}.}
          \resumeItemListEnd
      \resumeProjectHeading
          {\textbf{Meeting transcription using distributed microphones}}{April 2019 -- Present}
          \resumeItemListStart
            \resumeItem{Developed a meeting transcription framework based on distributed microphones, e.g., smartphones or tablets \cite{horiguchi2020utterancewise}.}
            \resumeItem{Developed its component technologies including speech separation \cite{horiguchi2021blockonline} and speaker diarization \cite{horiguchi2022multichannel,horiguchi2021towards,horiguchi2020endtoend,horiguchi2021endtoend}.}
            \resumeItem{Joined the Seventh Frederick Jelinek Memorial Summer Workshop (JSALT) as a part-time member, especially for discussion and technical support about speaker diarization \cite{horiguchi2021endtoend}.}
            \resumeItem{Won the second prize in the Third DIHARD Speech Diarization Challenge (DIHARD III).}
          \resumeItemListEnd
      \resumeProjectHeading
          {\textbf{Multi-modal recognition and human robot interaction}}{April 2017 -- March 2019}
          \resumeItemListStart
            \resumeItem{Developed multi-modal (especially audio-visual) functions for Hitachi's humanoid robot named EMIEW3, including multi-modal response obligation detection \cite{horiguchi2019multimodal} and cross-modal biometric search between faces and voices \cite{horiguchi2018facevoice}.}
          \resumeItemListEnd
      \resumeProjectHeading
          {\textbf{Personalized image recognition}}{April 2014 -- March 2017}
          \resumeItemListStart
            \resumeItem{Developed personalized image classifiers, especially for daily life-logging service \cite{horiguchi2018personalized}.}
            \resumeItem{Investigated better image representations for more general settings using convolutional neural networks \cite{horiguchi2020significance}.}
          \resumeItemListEnd
    \resumeSubHeadingListEnd
\end{comment}

\section{Publications}
  \resumeSubHeadingListStart
    \resumeShortSubheading
      {Journal Articles (first author)}{}
      \nocite{*}
      %\printbibliography
      \printbibliography[type=article,heading=none,keyword=first]
    \resumeShortSubheading
      {Journal Articles (co-author)}{}
      \printbibliography[type=article,heading=none,notkeyword=first]
    \resumeShortSubheading
      {Conference \& Workshop Proceedings (Peer-reviewed, first author)}{}
      \printbibliography[type=inproceedings,heading=none,keyword=first]
    \resumeShortSubheading
      {Conference \& Workshop Proceedings (Peer-reviewed, co-author)}{}
      \printbibliography[type=inproceedings,heading=none,notkeyword=first]
    \resumeShortSubheading
      {Preprints}{}
      \printbibliography[type=misc,heading=none]
    \resumeShortSubheading
      {\textbf{Invited Talk}}{}
      \resumeItemListStart
        \resumeItem{``Speaker Diarization: A Key to Solving Cocktail Party Problem,'' in IEEE RO-MAN Workshop, Aug. 28, 2023.}
        \resumeItem{``Face-Voice Matching Using Cross-Modal Embeddings,'' in MIRU, July 29--Aug. 1, 2020 (in Japanese).}
      \resumeItemListEnd
  \resumeSubHeadingListEnd

\section{Awards}
  \resumeItemListStart
    \resumeItem{IEEE SPS Young Author Best Paper Award, 2025. (Related to \cite{horiguchi2022encoderdecoder})}
    \resumeItem{IEEE SPS Japan Young Author Best Paper Award, 2025. (Related to \cite{horiguchi2023online})}
    \resumeItem{Honorable Mention Award at IEEE Spoken Language Technology Workshop, 2024. (Related to \cite{horiguchi2024recursive})}
    \resumeItem{Itakura Prize Innovative Young Researcher Award, The Acoustical Society of Japan, 2023. (Related to \cite{horiguchi2022encoderdecoder})}
    \resumeItem{ITE Outstanding Research Presentation Award, The Institute of Image Information and Television Engineers, 2017. (Related to \cite{horiguchi2020significance})}
  \resumeItemListEnd

\section{Competitions}
  \resumeItemListStart
    \resumeItem{2nd prize in the 8th CHiME Speech Separation and Recognition Challenge (CHiME-8) Task 1, 2024.}
    \resumeItem{2nd prize in the Third DIHARD Speech Diarization Challenge (DIHARD III), 2021  (as a lead author).}
    \resumeItem{2nd prize in the 5th CHiME Speech Separation and Recognition Challenge (CHiME-5), 2018.}
  \resumeItemListEnd

%\section{Patents and Applications}
%  \resumeItemListStart
%    \resumeItem{US Patent 11,107,476 B2 ``Speaker estimation method and speaker estimation device,'' August 31, 2021.}
%    \resumeItem{Contributed to six patent applications at Hitachi.}
%  \resumeItemListEnd
  
\section{Professional Activities/Services}
  \resumeSubHeadingListStart
    \resumeShortSubheading{Board member}{}{}
      \resumeItemListStart
        \resumeItem{IPSJ SIG-SLP (2025.4--Present)}
        \resumeItem{ASJ-SP (2025.4--Present)}
        \resumeItem{IEICE Technical Committee on Speech (2025.6--Present)}
      \resumeItemListEnd
    \resumeShortSubheading{Organizer}{}{}
      \resumeItemListStart
        \resumeItem{The Joint Workshop on HSCMA and CHiME 2026 (a satellite workshop of ICASSP 2026)}
      \resumeItemListEnd
    \resumeShortSubheading{Session chair}{}{}
      \resumeItemListStart
        \resumeItem{ICASSP (2022, 2025)}
        \resumeItem{Interspeech 2025)}
      \resumeItemListEnd
  \resumeSubHeadingListEnd
    
\section{Review Experiences}
 \begin{itemize}[leftmargin=0.15in, label={}]
    \small{\item{
     \textbf{Journals}{: IEEE/ACM Transactions on Audio, Speech, and Language processing, IEEE Transactions on Multimedia, IEEE Transactions on Neural Networks and Learning Systems, IEEE Transactions on Pattern Recognition and Machine Intelligence, IEEE Transactions on Affective Computing, Computer Speech \& Language, Speech Communication, EURASIP Journal on Audio, Speech and Language Processing, Neural Networks} \\
     \textbf{Conferences/Workshops}{: ICASSP, EUSIPCO, ASRU, SLT, WASPAA, DCASE, MLSP, MMSP, APSIPA ASC, ACMMM, RO-MAN, ICML}
    }}
 \end{itemize}
%
%-----------PROGRAMMING SKILLS-----------
%\section{Technical Skills}
% \begin{itemize}[leftmargin=0.15in, label={}]
%    \small{\item{
%     \textbf{Languages}{: Python, C/C++, Bash, LaTeX, JavaScript, HTML, CSS} \\
%     \textbf{Frameworks}{: PyTorch, Chainer, Caffe, Kaldi} \\
%     \textbf{Software \& Tools}{: Git, Vim}
%    }}
% \end{itemize}

\section{Language Skills}
 \begin{itemize}[leftmargin=0.15in, label={}]
    \small{\item{Japanese (native), English (business)}}
 \end{itemize}

%-------------------------------------------
\end{document}
